\section*{الفصل \ref{ch:b3:stable-isotopes} --- جيوكيمياء النظائر المستقرّة}

\begin{solution}{exo:b3:stable-isotopes:1}
$\alpha = 1030.5/997.9 = 1.032669$; $1000\ln\alpha = 32.15$; $\Delta = 30.5 -
(-2.1) = \qty{32.6}{\permil}$. يختلفان بمقدار \qty{0.45}{\permil}، ويطابق
$(\delta_A^2 - \delta_B^2)/2000 = 0.46$ ما ورد في الاقتراح.
\end{solution}

\begin{solution}{exo:b3:stable-isotopes:2}
$\delta^{18}\mathrm O = 0.69\times(-25) - 13.6 = \qty{-30.9}{\permil}$;
$\delta\mathrm D = 5.6\times(-25) - 100 = \qty{-240}{\permil}$; $d = -240 +
8\times30.85 = \qty{6.8}{\permil}$.
\end{solution}

\begin{solution}{exo:b3:stable-isotopes:3}
$\Delta_B = -26.0 - 1.0 = \qty{-27.0}{\permil}$; $f_{\mathrm{org}} = (-5 -
1.0)/(-27.0) = 0.22$.
\end{solution}

\begin{solution}{exo:b3:stable-isotopes:4}
$T^2 = 0.0391\times10^6/(0.600 - 0.154) = \num{87668}$, $T = \qty{296.1}{K}$,
\qty{22.9}{\celsius}.
\end{solution}

\begin{solution}{exo:b3:stable-isotopes:5}
$f^{0.01} = 967/987$، $f = \exp(100\ln0.97974) = 0.129$: \qty{87}{\percent} من
البخار قد تساقط منه. المطر: $1.0100\times967 - 1000 =
\qty{-23.3}{\permil}$.
\end{solution}

\begin{solution}{exo:b3:stable-isotopes:6}
$\alpha = 1000/960 = 1.041667$; $1000\ln\alpha = 40.82$; $1000(\alpha - 1) =
41.67$؛ $\Delta = 40$. يخطئ كلا التقريبين بنحو \qty{0.8}{\permil}،
والثاني بمقدار ضعف الأول: عند \qty{40}{\permil} يصبح التوسّع إلى
الرتبة الأولى غير كافٍ.
\end{solution}

\begin{solution}{exo:b3:stable-isotopes:7}
$\delta_c = (1.50 + 29.99)/0.97001 = \qty{32.46}{\permil}$ VSMOW؛ $1000\ln\alpha
= 1000\ln1.032464 = 31.95$; $T = 18030/(31.95 + 32.42) = \qty{280.1}{K}$,
\qty{7.0}{\celsius}. يتغيّر التجزّؤ بنحو \qty{0.23}{\permil} لكلّ
كلفن، فإنّ \qty{0.1}{\permil} في الماء يقابلها نحو \qty{0.4}{\celsius}.
\end{solution}

\begin{solution}{exo:b3:stable-isotopes:8}
حجم الجليد $125\times0.011 = \qty{1.375}{\permil}$؛ أمّا الباقي، $1.77 - 1.375 =
\qty{0.40}{\permil}$. قرب \qty{2}{\celsius} يبلغ الميل $18030/275.15^2 =
\qty{0.238}{\permil.K^{-1}}$، فإنّ الماء العميق كان
أبرد بنحو \qty{1.7}{\celsius}.
\end{solution}

\begin{solution}{exo:b3:stable-isotopes:9}
$\Delta^{33}\mathrm S = 8.2 - 0.515\times10.0 = \qty{3.05}{\permil}$، بعيدة عن 0:
الكبريت المستقل عن الكتلة، وهو لا يتكوّن إلا في غياب درع الأوزون. تكوّن البيريت
في غلاف جوي خالٍ من أكسجين معتدّ به، قبل نحو
\qty{2.45}{\Ga}.
\end{solution}

\begin{solution}{exo:b3:stable-isotopes:10}
يُحفَظ النظير الثقيل في النظام المغلق: $R_0 = fR_v + (1 -
f)\alpha R_v$، ومنه $R_v/R_0 = 1/[f + \alpha(1 - f)]$. من أجل $f = 0.2$، $\alpha =
1.01$: $1/1.008$، فيكون البخار أخفّ بمقدار \qty{7.9}{\permil}؛ ويعطي \omterm{prop:b3:stable-isotopes:rayleigh-distillation}{تقطير رايلي}
$0.2^{0.01}$، \qty{16.0}{\permil} أخفّ. وإزالة كلّ مكثَّف فورًا
يدفع البخار أبعد مرّتين.
\end{solution}

\begin{solution}{exo:b3:stable-isotopes:11}
$M = m(\delta_s - \Delta)/\Delta = 2000\times(-60 + 2.5)/(-2.5) = \qty{46000}{Gt}$.
ثمّ $m = 46000\times(-2.5)/(-25 + 2.5) = \qty{5100}{Gt}$، أقلّ قليلًا من
المنشورة من 6000 إلى \qty{12000}{Gt}، ويقابل طرفها الأعلى
الانحراف الأكبر.
\end{solution}

\begin{solution}{exo:b3:stable-isotopes:12}
$\delta_v = (1000 + \delta_{v,0})\exp\int(\alpha - 1)\,\mathrm d\ln f - 1000$، ومنه
$\mathrm d\delta_v = (1000 + \delta_v)(\alpha - 1)\,\mathrm d\ln f$، و$f =
e_s/e_{s,0}$ يعطي $\mathrm d\ln f = \mathrm d\ln e_s$. عند \qty{-20}{\celsius}،
$\alpha = 1.01872$: $0.01872\times950.9\times0.0960 =
\qty{1.71}{\permil.K^{-1}}$ للبخار.
\end{solution}

\begin{solution}{pb:b3:stable-isotopes:1}
\textbf{1.} $\alpha^{18}$: 1.01605, 1.01691, 1.01779, 1.01872, 1.01968,
1.02068; $\alpha^{\mathrm D}$: 1.1424, 1.1524, 1.1630, 1.1744, 1.1866, 1.1997
(عند $-5$ إلى \qty{-30}{\celsius}).

\textbf{2.} $980\times0.425^{0.01605} - 1000 = \qty{-33.4}{\permil}$.

\textbf{3.} $1.01691\times966.6 - 1000 = \qty{-17.0}{\permil}$.

\textbf{4.} البخار $966.6\times(0.3976)^{0.01779} - 1000 = \qty{-49.1}{\permil}$؛
الثلج $1.01872\times950.9 - 1000 = \qty{-31.3}{\permil}$.

\textbf{5.} البخار $950.9\times(0.3680)^{0.01968} - 1000 = \qty{-67.6}{\permil}$؛
الثلج $1.02068\times932.4 - 1000 = \qty{-48.3}{\permil}$.

\textbf{6.} $(932.4/980)\times0.0622 = 0.059$: \qty{5.9}{\percent} من الأكسجين-18،
مقابل \qty{6.2}{\percent} من الماء: فقد البخار
النظير الثقيل أسرع.

\textbf{7.} $(-17.0 + 48.3)/20 = \qty{1.57}{\permil}$ لكلّ درجة.

\textbf{8.} $t = (-48.3 + 13.6)/0.69 = \qty{-50}{\celsius}$. العلاقة صالحة
لدرجات حرارة السطح؛ فوق الغطاء الجليدي يكون السطح أبرد بكثير من
الهواء حيث يتكوّن الثلج، ويتجاهل النموذج أثر الخلط والتأثيرات الحركية.

\textbf{9.} البخار $-247.5$، $-352.5$، \qty{-462.7}{\permil}؛ الثلج $\delta\mathrm D
= -132.8$, $-239.6$, \qty{-355.4}{\permil}.

\textbf{10.} $d = -132.8 + 8\times17.02 = 3.3$؛ ثمّ 10.7
و\qty{31.3}{\permil}: يزداد الفائض نحو الداخل. ميل
التقطير المحلي، $\mathrm d\delta\mathrm D/\mathrm d\delta^{18}\mathrm O = (\alpha^{\mathrm D} -
1)(1000 + \delta\mathrm D)/[(\alpha^{18} - 1)(1000 + \delta^{18}\mathrm O)]$، ينخفض دون
8 مع ازدياد خفّة البخار كثيرًا ($1000 + \delta\mathrm D$ يتقلّص أسرع بكثير من
$1000 + \delta^{18}\mathrm O$)، فينحرف الثلج فوق الخط.

\textbf{11.} يعطي الخطّ $8\times(-48.34) + 10 = \qty{-376.7}{\permil}$؛
الثلج أعلى منه بمقدار \qty{21}{\permil}. يتضمّن النموذج التجزّؤ التوازني وحده،
أمّا الثلج الحقيقي فيتكوّن في هواء فوق التشبّع، بأثر حركي،
وتختلط الكتل الهوائية بغيرها في الطريق.

\textbf{12.} $\delta_c = (3.25 + 29.99)/0.97001 = 34.27$؛ $1000\ln\alpha =
33.69$؛ $T = 18030/66.11 = \qty{272.7}{K}$، \qty{-0.4}{\celsius}.

\textbf{13.} $\delta_c = 36.09$، $1000\ln\alpha = 35.46$، $T = \qty{265.6}{K}$،
\qty{-7.5}{\celsius}: أبرد ممّا تستطيع مياه البحر أن تبلغه دون أن تتجمّد. ويكون
التغيّر لا يمكن أن يكون كلّه بسبب درجة الحرارة.

\textbf{14.} $\delta_w = 125\times0.011 = \qty{+1.375}{\permil}$؛ $1000\ln(1036.09/1001.375)
= 34.08$؛ $T = 18030/66.50 = \qty{271.1}{K}$، \qty{-2.0}{\celsius}: تبريد بمقدار
\qty{1.6}{\celsius}.

\textbf{15.} $\Delta_{47} = 0.0391\times10^6/272.71^2 + 0.154 = 0.6797$
و$0.0391\times10^6/271.12^2 + 0.154 = 0.6859$: فرق قدره
\qty{0.006}{\permil}.

\textbf{16.} تقيس كيفية توزّع النظائر الثقيلة داخل
المعدن، وهو توازن داخلي يعتمد على درجة الحرارة وحدها.

\textbf{17.} $f_{\mathrm{org}} = (-5 - 0.5)/(-28) = 0.20$.

\textbf{18.} $m = 46000\times(-3.0)/(-60 + 3.0) = \qty{2400}{Gt}$.

\textbf{19.} $46000\times(-3.0)/(-25 + 3.0) = \qty{6300}{Gt}$.

\textbf{20.} $46000\times(-3.0)/(-5 + 3.0) = \qty{69000}{Gt}$، أكبر من
الخزّان نفسه، في زمن قصير جيولوجيًّا: غير معقول.

\textbf{21.} عمق ومدى ذوبان الكربونات في اللبّ، وهو ما
يقيس مقدار الكربون الذي دخل المحيط، بمعزل عن تركيبه؛
ومع حجم \omterm{def:b3:stable-isotopes:cie}{انحراف نظائر الكربون} يُحدَّد $\delta_s$.

\textbf{22.} الثلج عند دوم إيلسه يملك $\delta^{18}\mathrm O =
\textbf{\qty[text-series-to-math=true]{-48.3}{\permil}}$ (\qty{-48.4}{\permil}
مع تغيّر $\alpha$ تغيّرًا مستمرًّا).
\end{solution}
